\documentclass[10pt,a4paper]{amsart}
\usepackage{amsmath,amssymb,amsthm,graphicx}
% Retained minimal preamble for the Stacks Project chapter sites-cohomology.tex (Cohomology on Sites).
% Verbatim excerpt of the chapter preamble: Xy-pic setup (native preamble lines 14-18) and the native
% theorem declarations with the mathematical macros (native preamble lines 154-217). Copied mechanically;
% nothing added, removed or redefined. The wrapper class and the amsmath/amssymb/amsthm/graphicx base
% packages are supplied by the assembler, so the class/externaldocument part of the native preamble is omitted.
% For commutative diagrams we use Xy-pic
\usepackage[all]{xy}

% We use 2cell for 2-commutative diagrams.
\xyoption{2cell}
% Theorem environments.
%
\theoremstyle{plain}
\newtheorem{theorem}[subsection]{Theorem}
\newtheorem{proposition}[subsection]{Proposition}
\newtheorem{lemma}[subsection]{Lemma}

\theoremstyle{definition}
\newtheorem{definition}[subsection]{Definition}
\newtheorem{example}[subsection]{Example}
\newtheorem{exercise}[subsection]{Exercise}
\newtheorem{situation}[subsection]{Situation}

\theoremstyle{remark}
\newtheorem{remark}[subsection]{Remark}
\newtheorem{remarks}[subsection]{Remarks}

\numberwithin{equation}{subsection}

% Macros
%
\def\lim{\mathop{\mathrm{lim}}\nolimits}
\def\colim{\mathop{\mathrm{colim}}\nolimits}
\def\Spec{\mathop{\mathrm{Spec}}}
\def\Hom{\mathop{\mathrm{Hom}}\nolimits}
\def\Ext{\mathop{\mathrm{Ext}}\nolimits}
\def\SheafHom{\mathop{\mathcal{H}\!\mathit{om}}\nolimits}
\def\SheafExt{\mathop{\mathcal{E}\!\mathit{xt}}\nolimits}
\def\Sch{\mathit{Sch}}
\def\Mor{\mathop{\mathrm{Mor}}\nolimits}
\def\Ob{\mathop{\mathrm{Ob}}\nolimits}
\def\Sh{\mathop{\mathit{Sh}}\nolimits}
\def\NL{\mathop{N\!L}\nolimits}
\def\CH{\mathop{\mathrm{CH}}\nolimits}
\def\proetale{{pro\text{-}\acute{e}tale}}
\def\etale{{\acute{e}tale}}
\def\QCoh{\mathit{QCoh}}
\def\Ker{\mathop{\mathrm{Ker}}}
\def\Im{\mathop{\mathrm{Im}}}
\def\Coker{\mathop{\mathrm{Coker}}}
\def\Coim{\mathop{\mathrm{Coim}}}

% Boxtimes
%
\DeclareMathSymbol{\boxtimes}{\mathbin}{AMSa}{"02}

%
% Macros for moduli stacks/spaces
%
\def\QCohstack{\mathcal{QC}\!\mathit{oh}}
\def\Cohstack{\mathcal{C}\!\mathit{oh}}
\def\Spacesstack{\mathcal{S}\!\mathit{paces}}
\def\Quotfunctor{\mathrm{Quot}}
\def\Hilbfunctor{\mathrm{Hilb}}
\def\Curvesstack{\mathcal{C}\!\mathit{urves}}
\def\Polarizedstack{\mathcal{P}\!\mathit{olarized}}
\def\Complexesstack{\mathcal{C}\!\mathit{omplexes}}
% \Pic is the operator that assigns to X its picard group, usage \Pic(X)
% \Picardstack_{X/B} denotes the Picard stack of X over B
% \Picardfunctor_{X/B} denotes the Picard functor of X over B
\def\Pic{\mathop{\mathrm{Pic}}\nolimits}
\def\Picardstack{\mathcal{P}\!\mathit{ic}}
\def\Picardfunctor{\mathrm{Pic}}
\def\Deformationcategory{\mathcal{D}\!\mathit{ef}}
\usepackage{cleveref}
\crefname{originallemma}{original lemma}{original lemmas}
\Crefname{originallemma}{Original lemma}{Original lemmas}
\makeatletter
\newlabel{lemma-sheafification-cohomology}{{0BKV}{source}{}{sourcepointer-1}{}}
\newlabel{lemma-sheafification-cohomology@cref}{{[originallemma][0][]0BKV}{[1][0][]source}{}{}{}}
\makeatother
\crefname{originallemma}{original lemma}{original lemmas}
\Crefname{originallemma}{Original lemma}{Original lemmas}
\makeatletter
\newlabel{lemma-RGamma-commutes-with-Rlim}{{0D6K}{source}{}{sourcepointer-2}{}}
\newlabel{lemma-RGamma-commutes-with-Rlim@cref}{{[originallemma][0][]0D6K}{[1][0][]source}{}{}{}}
\makeatother
\title{Injectivity of sections of a derived limit on a suitable object}
\author{The Stacks Project Authors}
\date{stacks-project:sites-cohomology.tex:local-snapshot-20261011; GNU FDL 1.2 or later}
\begin{document}
\maketitle
\noindent\textit{Editorial note.} All delivered mathematical text is the author's own, copied from the pinned Stacks Project chapter file sites-cohomology.tex; no statement or proof step was rewritten, reordered or removed.\par\medskip
\noindent\textit{Editorial note.} The proof is the author's complete argument; both hypotheses are used explicitly and the conclusion is read off the exact sequence.\par\medskip
\noindent\textit{Editorial note.} Every labelled reference occurring inside the counted spans is carried in assembly.external\_references as a non-load-bearing navigation pointer (same-chapter ones at their definition site, cross-chapter ones at the invocation inside the counted span), so no load-bearing mathematics is replaced by a stub. No bibliography citation occurs inside the counted spans.\par\medskip
\noindent\textit{Editorial note.} The checked chapter list for same-name collisions is the set of stacks chapters staged in this repository: chow, cohomology, divisors, etale-cohomology, flat, more-algebra, more-morphisms, morphisms, spaces-more-morphisms, varieties together with this card's chapter sites-cohomology.\par\medskip
\noindent\textit{Editorial note.} The preamble is the chapter's own retained minimal preamble (sources/stacks-sites-cohomology/retained-preamble.tex), a verbatim excerpt of the native preamble: the Xy-pic setup and the native theorem declarations with the mathematical macros.\par\medskip
\noindent\textit{Allowed standard prerequisites (editorial).}\begin{itemize}
\item Derived limit of an inverse system of objects of the derived category and its short exact sequence. Same chapter together with the chapter Derived Categories, the defining triangle 0 -> R\textasciicircum{}1 lim -> R lim -> lim and its long exact cohomology sequence. Supplies the exact sequence from which the injectivity is read off.
\item Vanishing of R\textasciicircum{}1 lim for the cohomology of the objects of the system. Same chapter, the criterion for the vanishing of R\textasciicircum{}1 lim of an inverse system of abelian groups, which is hypothesis (1) of the statement. Supplies the vanishing that makes the relevant term disappear.
\end{itemize}\medskip
\ifdefined\hypertarget\hypertarget{sourcepointer-1}{}\fi
\noindent\textit{Source reference.} Original lemma 0BKV: sources/stacks-sites-cohomology/pinned-original-source.tex, source lines 4349--4349. Navigation pointer invoked at proof line 5296; the counted span contains the pointer only and the cited result is not reproduced inside it, so the pointer is registered as a non-load-bearing navigation stub in the established form for Stacks material. No bibliography citation is involved.\par\medskip
\ifdefined\hypertarget\hypertarget{sourcepointer-2}{}\fi
\noindent\textit{Source reference.} Original lemma 0D6K: sources/stacks-sites-cohomology/pinned-original-source.tex, source lines 5126--5126. Navigation pointer invoked at proof line 5238; the counted span contains the pointer only and the cited result is not reproduced inside it, so the pointer is registered as a non-load-bearing navigation stub in the established form for Stacks material. No bibliography citation is involved.\par\medskip
\begin{lemma}
\label{lemma-cohomology-derived-limit-injective}
Let $(\mathcal{C}, \mathcal{O})$ be a ringed site. Let $(K_n)$ be an
inverse system in $D(\mathcal{O})$. Let $V \in \Ob(\mathcal{C})$
and $m \in \mathbf{Z}$. Assume there exist an integer $n(V)$
and a cofinal system $\text{Cov}_V$ of coverings of $V$ such that
for $\{V_i \to V\} \in \text{Cov}_V$
\begin{enumerate}
\item $R^1\lim H^{m - 1}(V_i, K_n) = 0$, and
\item $H^m(V_i, K_n) \to H^m(V_i, K_{n(V)})$ is injective
for $n \geq n(V)$.
\end{enumerate}
Then the map on sections $H^m(R\lim K_n)(V) \to H^m(K_{n(V)})(V)$ is injective.
\end{lemma}

\begin{proof}
Let $\gamma \in H^m(R\lim K_n)(V)$ map to zero in
$H^m(K_{n(V)})(V)$. Since $H^m(R\lim K_n)$ is the sheafification of
$U \mapsto H^m(U, R\lim K_n)$ (by Lemma \ref{lemma-sheafification-cohomology})
we can choose $\{V_i \to V\} \in \text{Cov}_V$
and elements $\tilde\gamma_i \in H^m(V_i, R\lim K_n)$ mapping to
$\gamma|_{V_i}$.
Then $\tilde\gamma_i$ maps to $\tilde\gamma_{i, n(V)} \in H^m(V_i, K_{n(V)})$.
Using that $H^m(K_{n(V)})$ is the sheafification of
$U \mapsto H^m(U, K_{n(V)})$
(by Lemma \ref{lemma-sheafification-cohomology} again)
we see that after replacing $\{V_i \to V\}$ by a refinement
we may assume that $\tilde\gamma_{i, n(V)} = 0$ for all $i$.
For this covering we consider the short exact sequences
$$
0 \to
R^1\lim H^{m - 1}(V_i, K_n) \to H^m(V_i, R\lim K_n) \to
\lim H^m(V_i, K_n) \to 0
$$
of Lemma \ref{lemma-RGamma-commutes-with-Rlim}.
By assumption (1) the group on the left is zero and by
assumption (2) the group on the right maps injectively
into $H^m(V_i, K_{n(V)})$. We conclude $\tilde\gamma_i = 0$
and hence $\gamma = 0$ as desired.
\end{proof}

\end{document}
